\scanpage{69}
% letterhead
\noindent U N I V E R S I T E\hfill Montpellier, le 26 Juin 1978\\
DES SCIENCES ET TECHNIQUES\\
DU LANGUEDOC\\
\texttt{-{}-{}-{}-{}-{}-{}-{}-{}-{}-{}-{}-}\\
Service de la Scolarité\\
MR/GJ\\
Diffusion : Liste A\\
N°:4000

% printed form
\uline{AVIS DE SOUTENANCE D'UNE THESE DE DOCTORAT D'UNIVERSITE}

\uline{Nom et prénom du candidat} : PAYEN Jean-Luc

\uline{Titre de la thèse} : Contribution à l'étude de la protection des surfaces métalliques par les composés oxyfluorés du phosphore V.

\uline{Date, heure et lieu de soutenance prévus} :\\
13 JUILLET 1978 à 11 Heures\\
Ecole Nationale Supérieure de Chimie.

\uline{Nom du Directeur de thèse} :\\
Monsieur le Professeur L. COT.

\uline{Désignation du laboratoire où la thèse est préparée} :\\
Laboratoire de Chimie minérale appliquée - Chimie des Matériaux, \quad E. N. S. C. M.

\uline{Analyse de l'ouvrage} :

Les sels de l'anion $\mathrm{PO3F}^{2-}$ peuvent être utilisés lors de l'opération de post-traitement après phosphatation en substitution des sels de chrome VI hautement toxiques.

Une étude de la nature de la surface crée par phosphatation et passivation par $\mathrm{K2PO3F}$ a été entreprise par différentes techniques : ESCA,Microscopies électronique , et ionique. Elle montre une réorganisation physico chimique de la couche.

\uline{Durée approximative des recherches ayant abouti à la thèse} : 2 ans

\uline{Composition du jury proposé} :\\
MM.\qquad :\\
M.\quad MAURIN\\
L.\quad COT\\
P.\quad DESRE\\
J.\quad DURAND\\
X.\quad GRANDIN DE L'EPREVIER\\
B.\quad PARANT\\
P.R\quad PHILIP

Le Secrétaire Général,\\
\handwritten{\unsure{E Siau}}\\
\uline{E. SIAU}
% archive sheet number: 68
